\scanpage{106}
\mspage{- IV- \handwritten{878} -}
% archive sheet number: 105
\uline{un point de} $Y$ .

Il résulte en effet de (17.4.1,(b'')) qu'une $Y$-section $s$ de $X$ est une
immersion ouverte , et comme $X$ est un $Y$-schéma , $s$ est aussi une \deleted{\unsure{imme}}
immersion fermée (\uline{I},5.4.7) ; il en résulte que $s$ est un isomorphisme de
$Y$ sur un sous-préschéma de $X$ induit sur une partie ouverte et fermée de
$X$ , et comme \deleted{\unsure{$X$|$Y$}} $s(Y)$ est connexe , cést nécessairement une composante
connexe de $X$ . Le reste de la proposition est immédiat .

\handwritten{[}\uline{Proposition} (17.4.10) .- \uline{Soient} $g:Y\to S$ , $h:X\to S$ \uline{deux morphismes localement de présentation finie . Pour qu'un} $S$-\uline{morphisme} $f:X\to Y$ \uline{soit non ramifié
il faut et il suffit que pour tout} $s\in S$ , \uline{le morphisme} $f_s:h^{-1}(s)\to g^{-1}(s)$
\uline{déduit de $f$ par le} \deleted{\unsure{morph|chang}} \uline{changement de base} $\operatorname{Spec}(\kappa(s))\to S$ , \uline{soit non ramifié}.

Cela résulte aussitôt du fait que pour un morphisme localement de présentation finie , le fait d'être non ramifié est une propriété des \uline{fibres} de ce
morphisme\inserted{ (17.4.1,d))}, et \inserted{\handwritten{de ce}} que pour tout $y\in Y$ , on a $f^{-1}(y)=f_s^{-1}(y)$ si $s=g(y)$ .\handwritten{]}

\uline{17.5. Caractérisations des morphismes lisses} .

\uline{Théorème} (17.5.1) .- \uline{Soient} $f:X\to Y$ \uline{un morphisme localement de présentation finie , $x$ un point de} $X$ , $y=f(x)$ . \uline{Les conditions suivantes sont équivalentes} :

a) \uline{$f$ est lisse au point $x$} .

b) \uline{$f$ est plat au point $x$ et le} $\kappa(y)$-\uline{préschéma} $f^{-1}(y)$ \uline{est lisse sur} $\kappa(y)$
\uline{au point $x$} .

b') \uline{$f$ est régulier au point $x$} (6.8.1) .

c) \uline{L'anneau} $O_{X,x}$ \uline{est une} $O_{Y,y}$-\uline{algèbre formellement lisse pour les topologies discrètes} .

On peut se borner au cas où $Y=\operatorname{Spec}(A)$ , \deleted{\unsure{$X=\operatorname{Spec}(B)$}} $X=\operatorname{Spec}(C)$ , où $C=B/\underline J$
$B=A[T_1,..,T_n]$ étant une algèbre de polynômes et $\underline J$ un idéal de $B$ \uline{de type fini} . L'équivalence de a) et c) résulte alors de (\uline{0},22.6.4)\inserted{ \handwritten{a)$\Leftrightarrow$c)}} . D'autre
part , comme $f^{-1}(y)$ est un $\kappa(y)$-préschéma localement \deleted{noethérien}\inserted{\handwritten{de type fini}} , l'équivalence de b) et b') résulte de (6.8.6)\inserted{ \handwritten{a)$\Leftrightarrow$b)}} . Reste à voir que a) et b)
sont équivalentes .
\margin{\handwritten{Attention à l'\unsure{ambiguïté|ambiguité} dans ``lisse'' !}}
